% ============================================================
% 测试卷 — 第5章 三角函数
% ============================================================
\documentclass[12pt, a4paper]{ctexart}
\usepackage[margin=2cm]{geometry}
\usepackage{amsmath, amssymb}
\usepackage{enumitem}
\usepackage{fancyhdr}
\usepackage{xcolor}
\usepackage{tikz}
\usepackage{graphicx}
\usepackage{tabularx}
\usepackage{multicol}

\definecolor{mainblue}{RGB}{41, 98, 168}

\pagestyle{fancy}
\fancyhf{}
\fancyhead[L]{\small\textcolor{mainblue}{2025-2026学年度·第5章测试卷}}
\fancyhead[R]{\small\textcolor{mainblue}{数学（必修第一册）}}
\fancyfoot[C]{\small —\ \thepage\ —}
\renewcommand{\headrulewidth}{0.6pt}
\renewcommand{\headrule}{\hbox to\headwidth{\color{mainblue}\leaders\hrule height \headrulewidth\hfill}}

\setCJKmainfont{Hiragino Sans GB}
\setCJKsansfont{Heiti SC}

\newcommand{\blank}[1]{\underline{\hspace{#1}}}

\begin{document}

\begin{center}
{\Large\bfseries 第5章\quad 三角函数\quad 测试卷}\\[6pt]
{\normalsize 数学（必修第一册）}\\[4pt]
{\small 考试时间：120分钟\quad 满分：150分}
\end{center}

\vspace{2pt}
\begin{center}
\begin{tabular}{|c|c|c|c|c|}
\hline
题号 & 一 & 二 & 三 & 总分 \\ \hline
得分 & & & & \\ \hline
\end{tabular}
\end{center}
\vspace{4pt}

\noindent\textbf{注意事项：}
\begin{enumerate}[leftmargin=2em]
  \item 答题前，请在答题卡上填写姓名、班级和考号。
  \item 选择题请用2B铅笔在答题卡上涂黑对应选项。
  \item 非选择题请用黑色签字笔在答题卡指定区域作答。
\end{enumerate}

\vspace{6pt}

% ============ 一、单项选择题 ============
\noindent\rule{\textwidth}{0.4pt}\\[2pt]
{\large\bfseries 一、单项选择题}（本大题共8小题，每小题5分，共40分。每小题只有一个选项符合题目要求。）

\vspace{6pt}

\begin{enumerate}
\item $\sin\dfrac{\pi}{6}$ 的值为（\quad）
  \begin{enumerate}[label=\Alph*., itemsep=0pt]
    \item $\dfrac{1}{2}$ \quad \item $\dfrac{\sqrt{2}}{2}$ \quad \item $\dfrac{\sqrt{3}}{2}$ \quad \item $1$
  \end{enumerate}

\item 将 $200°$ 化为弧度为（\quad）
  \begin{enumerate}[label=\Alph*., itemsep=0pt]
    \item $\dfrac{2\pi}{3}$ \quad \item $\dfrac{5\pi}{6}$ \quad \item $\dfrac{10\pi}{9}$ \quad \item $\dfrac{7\pi}{6}$
  \end{enumerate}

\item 已知角 $\alpha$ 的终边经过点 $P(-3,4)$，则 $\sin\alpha$ 的值为（\quad）
  \begin{enumerate}[label=\Alph*., itemsep=0pt]
    \item $\dfrac{3}{5}$ \quad \item $-\dfrac{3}{5}$ \quad \item $\dfrac{4}{5}$ \quad \item $-\dfrac{4}{5}$
  \end{enumerate}

\item 函数 $y=2\sin x$ 的最大值和最小值分别为（\quad）
  \begin{enumerate}[label=\Alph*., itemsep=0pt]
    \item $1,-1$ \quad \item $2,-2$ \quad \item $2,0$ \quad \item $1,-2$
  \end{enumerate}

\item 函数 $y=\sin 2x$ 的最小正周期为（\quad）
  \begin{enumerate}[label=\Alph*., itemsep=0pt]
    \item $\pi$ \quad \item $2\pi$ \quad \item $\dfrac{\pi}{2}$ \quad \item $4\pi$
  \end{enumerate}

\item $\cos(\pi-\alpha)$ 等于（\quad）
  \begin{enumerate}[label=\Alph*., itemsep=0pt]
    \item $\cos\alpha$ \quad \item $-\cos\alpha$ \quad \item $\sin\alpha$ \quad \item $-\sin\alpha$
  \end{enumerate}

\item 已知 $\tan\alpha=2$，则 $\dfrac{\sin\alpha+\cos\alpha}{\sin\alpha-\cos\alpha}$ 的值为（\quad）
  \begin{enumerate}[label=\Alph*., itemsep=0pt]
    \item $3$ \quad \item $-3$ \quad \item $\dfrac{1}{3}$ \quad \item $-\dfrac{1}{3}$
  \end{enumerate}

\item 函数 $y=\sin\!\left(2x+\dfrac{\pi}{3}\right)$ 的图象可由 $y=\sin 2x$ 的图象（\quad）得到。
  \begin{enumerate}[label=\Alph*., itemsep=0pt]
    \item 向左平移 $\dfrac{\pi}{3}$ 个单位 \quad \item 向右平移 $\dfrac{\pi}{3}$ 个单位
    \item 向左平移 $\dfrac{\pi}{6}$ 个单位 \quad \item 向右平移 $\dfrac{\pi}{6}$ 个单位
  \end{enumerate}
\end{enumerate}

\vspace{6pt}

% ============ 二、多项选择题 ============
\noindent\rule{\textwidth}{0.4pt}\\[2pt]
{\large\bfseries 二、多项选择题}（本大题共4小题，每小题5分，共20分。全部选对得5分，部分选对得2分，有选错的得0分。）

\vspace{6pt}

\begin{enumerate}[resume]
\item 下列关于角的说法正确的是（\quad）
  \begin{enumerate}[label=\Alph*., itemsep=0pt]
    \item 第一象限角一定是锐角
    \item 钝角是第二象限角
    \item $-150°$ 是第三象限角
    \item $390°$ 与 $30°$ 的终边相同
  \end{enumerate}

\item 下列等式成立的是（\quad）
  \begin{enumerate}[label=\Alph*., itemsep=0pt]
    \item $\sin^2\alpha+\cos^2\alpha=1$
    \item $\dfrac{\sin\alpha}{\cos\alpha}=\tan\alpha$
    \item $\sin(\pi-\alpha)=\sin\alpha$
    \item $\cos(-\alpha)=-\cos\alpha$
  \end{enumerate}

\item 关于函数 $y=\cos x$，下列说法正确的是（\quad）
  \begin{enumerate}[label=\Alph*., itemsep=0pt]
    \item 最小正周期为 $2\pi$
    \item 在 $[0,\pi]$ 上单调递减
    \item 是偶函数
    \item 最大值为 $1$，最小值为 $-1$
  \end{enumerate}

\item 已知函数 $f(x)=\sin\!\left(x+\dfrac{\pi}{4}\right)$，则（\quad）
  \begin{enumerate}[label=\Alph*., itemsep=0pt]
    \item $f(x)$ 的最小正周期为 $2\pi$
    \item $f\!\left(\dfrac{\pi}{4}\right)=1$
    \item $f(x)$ 的图象关于 $x=\dfrac{\pi}{4}$ 对称
    \item $f(x)$ 在 $\left[-\dfrac{3\pi}{4},\dfrac{\pi}{4}\right]$ 上单调递增
  \end{enumerate}
\end{enumerate}

\vspace{6pt}

% ============ 三、填空题 ============
\noindent\rule{\textwidth}{0.4pt}\\[2pt]
{\large\bfseries 三、填空题}（本大题共4小题，每小题5分，共20分。）

\vspace{6pt}

\begin{enumerate}[resume]
\item $\sin 30°+\cos 60°=$ \underline{\hspace{2cm}}。

\item 已知 $\sin\alpha=\dfrac{3}{5}$，且 $\alpha$ 为第一象限角，则 $\cos\alpha=$ \underline{\hspace{2cm}}。

\item 函数 $y=3\sin\!\left(2x-\dfrac{\pi}{6}\right)$ 的振幅为 \underline{\hspace{2cm}}，最小正周期为 \underline{\hspace{2cm}}。

\item 若 $\sin\alpha+\cos\alpha=\dfrac{7}{5}$，则 $\sin\alpha\cdot\cos\alpha=$ \underline{\hspace{2cm}}。
\end{enumerate}

\vspace{6pt}

% ============ 四、解答题 ============
\noindent\rule{\textwidth}{0.4pt}\\[2pt]
{\large\bfseries 四、解答题}（本大题共6小题，共70分。解答应写出文字说明、证明过程或演算步骤。）

\vspace{6pt}

\begin{enumerate}[resume]
\item[17.]（本小题10分）\\
已知 $\tan\alpha=3$，求下列各式的值：
\begin{enumerate}[(1)]
  \item $\dfrac{2\sin\alpha+\cos\alpha}{\sin\alpha-\cos\alpha}$；
  \item $\sin^2\alpha+\sin\alpha\cos\alpha$。
\end{enumerate}

\vspace{6cm}

\item[18.]（本小题12分）\\
已知角 $\alpha$ 的终边经过点 $P(-1,\sqrt{3})$。
\begin{enumerate}[(1)]
  \item 求 $\sin\alpha$、$\cos\alpha$、$\tan\alpha$ 的值；
  \item 求 $\dfrac{\sin(\pi+\alpha)\cos(2\pi-\alpha)}{\tan(-\alpha)\sin(-\pi-\alpha)}$ 的值。
\end{enumerate}

\vspace{6cm}

\item[19.]（本小题12分）\\
求函数 $y=2\sin\!\left(2x+\dfrac{\pi}{3}\right)$ 的：
\begin{enumerate}[(1)]
  \item 最小正周期；
  \item 单调递增区间；
  \item 当 $x\in\left[0,\dfrac{\pi}{2}\right]$ 时的最大值和最小值。
\end{enumerate}

\vspace{6cm}

\item[20.]（本小题12分）\\
已知 $\sin\alpha=\dfrac{4}{5}$，$\alpha\in\left(\dfrac{\pi}{2},\pi\right)$。
\begin{enumerate}[(1)]
  \item 求 $\cos\alpha$ 和 $\tan\alpha$ 的值；
  \item 求 $\sin\!\left(\alpha+\dfrac{\pi}{3}\right)$ 的值。
\end{enumerate}

\vspace{6cm}

\item[21.]（本小题12分）\\
已知函数 $f(x)=A\sin(\omega x+\varphi)$（$A>0$，$\omega>0$，$|\varphi|<\dfrac{\pi}{2}$）的部分图象如图所示，其中最高点的纵坐标为 $2$，相邻的最高点与最低点之间的水平距离为 $\dfrac{\pi}{2}$，图象过点 $(0,1)$。
\begin{enumerate}[(1)]
  \item 求 $A$、$\omega$、$\varphi$ 的值；
  \item 求 $f(x)$ 的单调递减区间。
\end{enumerate}

\vspace{7cm}

\item[22.]（本小题12分）\\
已知函数 $f(x)=\sin^2 x+\sqrt{3}\sin x\cos x+\dfrac{1}{2}$。
\begin{enumerate}[(1)]
  \item 将 $f(x)$ 化为 $A\sin(\omega x+\varphi)+B$ 的形式（其中 $A>0$，$\omega>0$，$0<\varphi<\pi$）；
  \item 求 $f(x)$ 的最小正周期及最大值；
  \item 求 $f(x)$ 在区间 $\left[0,\dfrac{\pi}{2}\right]$ 上的值域。
\end{enumerate}

\vspace{7cm}
\end{enumerate}

\vfill
\begin{center}
\small\textcolor{gray}{本试卷仅供教学参考}
\end{center}

\end{document}
